\section{Introduction}
MILP has become a standard tool for automated cryptanalysis of block ciphers, especially when searching for optimal differential and related-key characteristics. In this setting, the choice of MILP solver often determines whether an analysis is practical or infeasible, since solver runtime can vary by orders of magnitude on the same model.

This paper focuses on the performance comparison of five MILP solvers on related-key differential models for the lightweight cipher ITUbee. We benchmark three open-source solvers (GLPK, HiGHS, SCIP) and two commercial solvers (Gurobi, CPLEX) on 8-, 10-, and 12-round instances. The comparison highlights how algorithmic design, preprocessing, and search heuristics affect solver efficiency on cryptanalytic MILP workloads.

ITUbee is a software-oriented lightweight Feistel cipher operating on 80-bit blocks and an 80-bit key \cite{karakoc2013itubee}. Its round structure includes key whitening, substitution, and linear diffusion layers, making it a suitable target for MILP-based related-key differential analysis. This paper uses ITUbee as a representative case study for solver benchmarking rather than as a proposal for a new cipher design.

The main contributions of this work are:
\begin{itemize}
    \item a direct empirical comparison of GLPK, HiGHS, SCIP, Gurobi, and CPLEX on a series of ITUbee-related MILP models;
    \item an analysis of how solver performance scales as the model difficulty increases from 8 to 12 rounds;
    \item practical recommendations for choosing a solver for MILP-based related-key differential cryptanalysis.
\end{itemize}

The results are intended to guide cryptanalysts who need to select the most effective solver for MILP-based analysis, with a clear emphasis on performance rather than cryptographic design.

\textbf{Organization of the paper}. Section 2 introduces MILP-based cryptanalysis and related-key differential attacks. Section 3 describes the ITUbee cipher and the MILP model. Section 4 presents the benchmark results and analysis. Section 5 concludes with practical solver recommendations.
